\chapter{Polycystic ovary syndrome (PCOS)}
\index{Polycystic ovary syndrome (PCOS)}

\medskip
\noindent\textit{Educational information; seek professional advice for individual care.}

\section*{What it is}
Polycystic ovary syndrome is a hormonal and metabolic condition associated with irregular ovulation, androgen effects, and sometimes polycystic-appearing ovaries. The name is misleading: ovarian cysts are not required for diagnosis.

\section*{Symptoms}
Possible features include irregular or absent periods, acne, excess facial or body hair, scalp hair thinning, difficulty conceiving, weight changes, and darkened or thickened skin in some people. Symptoms vary, and other conditions can cause similar changes.

\section*{Assessment}
Assessment includes menstrual and medical history, examination, pregnancy testing when relevant, selected hormone tests, and screening for related metabolic risks such as high blood pressure, abnormal glucose, and sleep apnoea. Ultrasound may be useful but is not always required. Diagnosis requires excluding other explanations.

\section*{Management}
Treatment is based on symptoms, health risks, and pregnancy goals. It may include sustainable activity and nutrition support, menstrual regulation, treatment for acne or excess hair, medicines for metabolic risk, and fertility treatment. Emotional wellbeing and disordered eating should be discussed sensitively.

\section*{When to seek medical care}
Arrange review for prolonged absence of periods, very heavy bleeding, new rapid androgen-related changes, infertility concerns, or distress. Seek urgent care for severe pelvic pain, fainting, heavy uncontrolled bleeding, or rapidly worsening illness.

\section*{Sources}
\begin{itemize}
\item American College of Obstetricians and Gynecologists: \url{https://www.acog.org/womens-health/faqs/polycystic-ovary-syndrome-pcos}
\end{itemize}
